\documentclass{article}

\usepackage[margin=2.5cm]{geometry}
\usepackage{times}
\usepackage{amsmath,amssymb,amsfonts}
\usepackage{mathtools}
\usepackage{booktabs}
\usepackage{graphicx}
\usepackage{xcolor}
\usepackage{hyperref}
\usepackage{microtype}
\usepackage{caption}
\usepackage{subcaption}
\usepackage{array}
\usepackage{multirow}
\usepackage{bbm}

\hypersetup{
  colorlinks=true,
  linkcolor=black,
  citecolor=black,
  urlcolor=blue
}

% Fix citation superscript — avoids font corruption with times
\newcommand{\citenum}[1]{%
  \textsuperscript{\rmfamily\upshape[#1]}}

% Figures uploaded as PDF to Overleaf root directory
\graphicspath{{./}}

\title{\textbf{A Wave-Cache Architecture with Phase-Aware Training\\
for Language and Genomic Sequence Modeling}}

\author{
  Yasong Fan\\
  Independent Researcher\\
  \texttt{arXiv: 2604.07716} $\cdot$ \texttt{github.com/YasongFan/FDM}
}

\date{}

\begin{document}

\maketitle

%------------------------------------------------------------------
\begin{abstract}
Sequence modeling architectures face a fundamental tension between
global state propagation and selective local retrieval, yet most
approaches couple these functions within a single operator.
We introduce the Fan Duality Model (FDM), which explicitly separates
these functions into a \emph{wave channel}---implementing
norm-preserving recurrent propagation via Givens-rotation phase
operators---and a \emph{cache channel}, implementing sparse associative
retrieval over a fixed set of $W$ local and $K$ global slots independent
of sequence length $N$.
This separation yields a fixed-capacity historical state and measured
runtime memory saturation beyond moderate context lengths, in contrast
to the linear theoretical growth of Transformer KV-cache requirements
(measured incremental memory saturates beyond moderate context lengths).

A central empirical finding is that jointly training these two channels
produces a \emph{gradient sink}: scan parameters dominate gradient flow,
leaving the cache undertrained.
We introduce \emph{phase-aware staged training}, in which the wave
channel is frozen during an initial phase to allow the cache to
specialise, with automatic transition triggered by a gradient-ratio
criterion.
This strategy substantially improves performance and indicates functional
specialisation between channels: under our sequence-level mutual
information estimator, the wave channel carries substantially more
information than the cache channel with respect to the evaluated
statistic, with cache attention concentrated in local distances
(39\% within distance 16 tokens).

On WikiText-103\citenum{1}, FDM achieves validation perplexity of 33.75
under phase-aware training, below a similarly sized Transformer reference
under the present training setup ($\sim$36--38).
Pretrained on human hg38 (2.94\,Gbp), FDM achieves $0.945 \pm 0.001$
on promoter recognition, $0.905 \pm 0.012$ and $0.775 \pm 0.002$ on
two enhancer benchmarks, and $0.811 \pm 0.009$ on OCR prediction,
compared with HyenaDNA\citenum{2} references of 0.856, 0.706, 0.711,
and 0.806, respectively (three random seeds).
Fine-tuning on ClinVar variant effect prediction yields
AUC $= 0.861 \pm 0.011$ (zero-shot: 0.51; CADD: 0.869).
Cross-species transfer to \emph{Drosophila} enhancer recognition yields
AUC $= 0.792$, compared to 0.736 for a CNN trained from scratch.
\end{abstract}

%------------------------------------------------------------------
\section{Introduction}

Biological sequence modeling---and sequence modeling more broadly---requires
integrating two distinct computational capacities: the propagation of
distributional patterns across long distances, and the precise retrieval
of specific past tokens on demand.
Standard Transformers\citenum{3} address both through full self-attention,
but the resulting linear growth in inference memory limits deployment on
long sequences.
Linear recurrent models such as S4\citenum{4}, Mamba\citenum{5}, and
RWKV\citenum{6} achieve fixed-size inference memory through state
compression, yet this compression involves an inherent structural tension:
a single linear operator cannot simultaneously act as a norm-preserving
global integrator (retaining all historical signal) and a selective filter
(discarding irrelevant content on demand).
We frame this as a tension between norm-preserving global propagation and
selective retention of sparse facts, and use it as the motivating
observation for architectural separation.

FDM resolves this structural tension by architectural separation.
The wave channel handles norm-preserving propagation; the cache channel
handles selective retrieval.
This decomposition is not merely architectural---jointly training these
channels produces a gradient sink effect in which scan parameters
dominate gradient flow, leaving the cache severely undertrained.
We introduce a phase-aware staged training strategy that freezes the
wave channel during an initial training phase, forcing gradient signal
through the cache circuit before joint optimisation resumes.

In language modeling, FDM achieves perplexity of 33.75 on
WikiText-103\citenum{1}, below a similarly sized Transformer reference
($\sim$36--38), with measured runtime memory saturating at moderate
context lengths.
Mutual information analysis indicates that the wave channel encodes
global sequence-level information while the cache channel develops
locality-biased retrieval patterns, consistent with the intended
functional separation.

We further apply FDM to genomic sequence modeling, where the
fixed-capacity memory characteristic allows processing of longer
regulatory sequences on a single GPU.
A model pretrained on 2.94\,Gbp of human reference genome achieves
consistent improvement over HyenaDNA\citenum{2} across all four genomic
benchmarks evaluated (mean AUC +0.089 on promoter recognition, +0.199
on enhancer classification).
Fine-tuning on ClinVar variant effect data approaches CADD, a tool
built on manually curated features.
Cross-species transfer experiments indicate partial generalisation to
\emph{Drosophila} regulatory elements.

The main contributions are:
(1) the FDM architecture with fixed-capacity historical state, motivated by
the structural tension between norm-preserving propagation and selective
retention;
(2) phase-aware staged training, which addresses the gradient sink problem
and substantially improves performance and functional channel separation;
(3) systematic evaluation across language and genomic sequence modeling
tasks with clear characterisation of task boundaries.

%------------------------------------------------------------------
\section{Related Work}

\paragraph{Linear recurrent models.}
S4\citenum{4}, Mamba\citenum{5}, Mamba-2\citenum{7}, RWKV\citenum{6},
and RetNet\citenum{8} achieve fixed inference memory via recurrent state
compression.
FDM extends this family with explicit wave--cache separation.
Relative to Mamba's selective state space mechanism, FDM separates global
propagation from local retrieval architecturally rather than through
input-dependent gating of a single state.

\paragraph{KV cache compression.}
H2O\citenum{9} and StreamingLLM\citenum{10} apply post-hoc heuristics
to compress Transformer KV caches.
FDM's cache learns selection end-to-end via the $s_\mathrm{eff}$ signal.
Longformer\citenum{11} and BigBird\citenum{12} implement
local$+$global attention patterns in full attention; FDM implements
analogous patterns within a linear recurrent framework with
$O(W{+}K)$ memory.

\paragraph{DNA language models.}
HyenaDNA\citenum{2} applies a Hyena convolution operator to genomic
sequences at long sequence lengths.
Nucleotide Transformer\citenum{13} and DNABERT-2\citenum{14} apply
BERT-style pretraining to DNA.
Evo\citenum{15} scales Hyena-based pretraining to multi-billion parameter
models on prokaryotic genomes.
FDM differs in its explicit wave--cache decomposition and in operating on
eukaryotic human reference sequence.

\paragraph{Staged training.}
Curriculum learning\citenum{16} and progressive training strategies have
been applied in various domains.
Phase-aware training is most closely related to the observation that
induction heads form in a two-phase process\citenum{17}; we observe an
analogous phenomenon in FDM.

%------------------------------------------------------------------
\section{Method}

\subsection{Architecture Overview}

An overview of the architecture and training procedure is shown in
Figure~\ref{fig:arch}.

\begin{figure}[t]
\centering
\includegraphics[width=\linewidth]{fig1_architecture}
\caption{
  \textbf{FDM architecture and phase-aware training.}
  Input tokens are processed by $L$ stacked layers, each comprising a
  wave channel (norm-preserving Givens-rotation scan) and a cache channel
  (sparse attention over $W$ local and $K$ global slots). Channel outputs
  are combined via a learned gate. Phase-aware training proceeds in two
  stages: Phase~0 freezes wave parameters, forcing gradient flow through
  the cache circuit; Phase~1 resumes joint optimisation after automatic
  transition at step~23{,}000.
}
\label{fig:arch}
\end{figure}

FDM processes tokens $x_1, \ldots, x_T$ through $L$ stacked layers,
each comprising a wave channel and a cache channel.
The wave channel output $h_t$ and cache channel output $c_t$ are combined
via a learned gate:
\begin{equation}
  x_t^{(l+1)} = h_t^{(l)} + g_t^{(l)} \odot c_t^{(l)},
  \quad
  g_t^{(l)} = \sigma\!\left(W_g h_t^{(l)}\right) \in (0,1)^D.
\end{equation}

\subsection{Wave Channel}

The wave channel implements the Fan Operator, a two-pass linear recurrent
scan with phase-preserving Givens rotations.
The first pass computes:
\begin{equation}
  h_t^{(1)} = (1 - p_t) \odot R(\theta_t) \cdot h_{t-1}
              + p_t \odot \bigl(W_r x_t + i\, W_i x_t\bigr),
\end{equation}
where $R(\theta_t) \in \mathbb{C}^{D \times D}$ is a unitary rotation
via $n_G$ Givens pairs satisfying $\|R(\theta)h\| = \|h\|$, and
$p_t \in (0, 0.5]$:
\begin{align}
  s_\mathrm{eff}(t) &= \beta_t \cdot \log(t+2)
    + W_\mathrm{pos}[\sin(\pi t/T),\, \cos(\pi t/T)]^\top, \\
  p_t &= \sigma\!\left(s_\mathrm{eff}(t) + \mu\right) \cdot 0.5
         + \varepsilon.
\end{align}
A second pass applies a perturbative correction (Born approximation)
sharing all weights with the first, yielding $h_t = h_t^{(2)} + \Delta_\mathrm{gate} \cdot x_t$;
full equations are given in Supplementary~S3.

\subsection{Cache Channel}

The cache channel applies sparse attention over $W$ local and $K$ global
positions:
\begin{align}
  \mathrm{mask}(i,j) &= \mathrm{local}(i,j) \vee \mathrm{global}(i,j),\\
  \mathrm{local}(i,j)
    &= \mathbbm{1}\!\left[\max(0,i{-}W) \le j < i\right], \\
  \mathrm{global}(i,j)
    &= \mathbbm{1}\!\left[j \in
       \mathrm{TopK}(\{s_\mathrm{eff}(k) : k < i\},\, K)\right].
\end{align}
Global slot selection reuses $s_\mathrm{eff}$ without additional
parameters.
Since $W$ and $K$ are constants independent of $N$, the cache size is
fixed at $O(W{+}K)$.
Language model: $W = 256$, $K = 16$; genomic model: $W = 128$, $K = 32$.

\subsection{Phase-Aware Staged Training}

Joint optimisation produces a gradient sink: $\Phi_\mathrm{wave}$
dominates gradient magnitude, starving $\Phi_\mathrm{cache}$.

\textbf{Phase~0.}
$\Phi_\mathrm{wave}$ is frozen; $\Phi_\mathrm{cache}$ is optimised at
$\eta_0 = 10^{-4}$.

\textbf{Automatic transition.}
\begin{equation}
  \|\nabla_\mathrm{embed}\| \;/\;
  \mathrm{mean}\!\left(
    \|\nabla_{W_q}\|, \|\nabla_{W_k}\|, \|\nabla_{W_v}\|\right)
  > \alpha \quad (\alpha = 1.2).
\end{equation}

\textbf{Phase~1.}
All parameters jointly optimised at $\eta_1 = 3 \times 10^{-5}$.
Transition occurs at step~23{,}000 in language modeling experiments;
PPL decreases from 36.8 to 33.75 over the remaining 17{,}000 steps.

\subsection{Genomic Pretraining and Fine-tuning}

Pretrained autoregressively on hg38 (chromosomes 1--7, 9--16, 18--22, X;
2.94\,Gbp; chromosome~8 excluded).
Vocabulary: \{A, T, G, C, N\}.
$d = 512$, 8~layers, 95{,}000 steps, sequence length~512, batch~8;
validation PPL $= 3.23$.
Fine-tuning: first~4 layers and embedding frozen; upper~4 layers and a
two-layer head optimised via AdamW ($\mathrm{lr} = 5 \times 10^{-5}$,
cosine schedule, patience~4).

%------------------------------------------------------------------
\section{Results}

\subsection{Language Modeling}

\begin{table}[h]
\centering
\caption{WikiText-103 validation perplexity. All FDM variants share
the same architecture; the Transformer reference uses matched training
data, sequence length, and optimiser.}
\label{tab:lm}
\begin{tabular}{lccc}
\toprule
Model & Params & Steps & Val PPL $\downarrow$ \\
\midrule
Transformer (reference) & $\sim$137M & 42K & $\sim$36--38 \\
FDM --- joint training  & 137M & 42K & 487.5 \\
FDM --- Freeze-Scan (Phase~0 only) & 137M & 44K & 64.9 \\
FDM --- phase-aware (full) & 137M & 40K & \textbf{33.75} \\
\bottomrule
\end{tabular}
\end{table}

\begin{figure}[h]
\centering
\includegraphics[width=\linewidth]{fig2_lm}
\caption{
  \textbf{Language modeling results on WikiText-103.}
  \textbf{(A)}~Training dynamics: joint baseline remains unstable
  throughout training; phase-aware FDM converges steadily, with the
  dashed line marking the Transformer reference ($\sim$36--38~PPL).
  The vertical dotted line marks the automatic phase transition at
  step~23{,}000.
  \textbf{(B)}~Validation PPL by training configuration. Phase-aware
  training reduces PPL from 487.5 to 33.75, below the Transformer
  reference. All variants share the same architecture; differences
  reflect optimisation strategy under the present setup.
}
\label{fig:lm}
\end{figure}

Joint training (PPL $= 487.5$) is the gradient sink baseline.
Phase~0 alone reaches 64.9; full phase-aware training reaches 33.75,
below the Transformer reference.
Transition at step~23{,}000 is automatic ($\alpha = 1.2$; PPL $= 36.8$
at transition).
The comparison primarily reflects differences in optimisation behaviour
and model design under the present setup.
Competitive SSM baselines trained on substantially larger corpora are not
directly comparable; see Section~\ref{sec:discussion}.

\begin{table}[h]
\centering
\caption{Freeze-Scan component ablation.}
\label{tab:ablation}
\begin{tabular}{lc}
\toprule
Configuration & Val PPL $\downarrow$ \\
\midrule
Baseline (joint)               & 487.5 \\
$+$ Freeze-Scan (Phase~0)      & 64.9  \\
$+$ $W_\mathrm{content}$ gate  & 60.8  \\
$+$ Recall loss                & 55.7  \\
$+$ Adaptive $\theta$          & 46.8  \\
$+$ Phase~1 joint              & \textbf{33.75} \\
\bottomrule
\end{tabular}
\end{table}

\begin{table}[h]
\centering
\caption{Runtime memory increment (MB) vs.\ sequence length.
FDM increment remains approximately constant above $N = 2{,}048$.}
\label{tab:memory}
\begin{tabular}{lcc}
\toprule
$N$ & TF KV cache (MB) & FDM increment (MB) \\
\midrule
2,048  & 56.6    & $+$396.9 \\
4,096  & 113.2   & $+$396.9 \\
16,384 & 453.0   & $+$396.9 \\
65,536 & 1,811.9 & $+$396.9 \\
\bottomrule
\end{tabular}
\end{table}

At $N = 65{,}536$, FDM requires $4.6\times$ less incremental memory.

\subsection{Wave--Cache Functional Analysis}

\begin{figure}[h]
\centering
\includegraphics[width=\linewidth]{fig3_analysis}
\caption{
  \textbf{Wave--cache functional analysis.}
  \textbf{(A)}~Mutual information between channel representations and a
  sequence-level surprisal statistic (MINE estimator). The wave channel
  carries substantially more information than the cache channel under
  this statistic.
  \textbf{(B)}~Cache attention mass by token distance, showing a
  locality bias (39\% within distance~16).
  \textbf{(C)}~Wave channel phase parameter $\theta$ (mean~$\pm$~std)
  across layers. Values cluster near $\theta \approx 0.50$
  ($\lambda \approx 12.5$ tokens) with no significant layer-to-layer
  variation, indicating operation at a single temporal scale rather than
  an emergent multi-scale hierarchy.
}
\label{fig:analysis}
\end{figure}

The wave channel carries $\mathrm{MI} = 0.474$ nats with sequence-level
surprisal statistics; the cache channel carries little measurable mutual
information under the same statistic.
Cache attention exhibits a locality bias: 39\% of attention mass falls
within distance~16 tokens, decreasing monotonically to 5.2\% at
128--256.
Wave channel $\theta$ values cluster near 0.50 across all layers
(effective wavelength $\lambda \approx 12.5$ tokens), indicating
operation at a consistent temporal scale.

\subsection{Genomic Benchmarks}

\begin{table}[h]
\centering
\caption{Genomic benchmark AUC (mean $\pm$ std, 3 seeds).
Tasks from \texttt{genomic\_benchmarks}\citenum{18}; protocol follows
HyenaDNA\citenum{2}.}
\label{tab:genomic}
\begin{tabular}{lcccc}
\toprule
Task & Seq.\ len & FDM-HG38 & HyenaDNA & $\Delta$ \\
\midrule
human\_nontata\_promoters & 251\,bp
  & $\mathbf{0.945 \pm 0.001}$ & 0.856 & $+$0.089 \\
human\_enhancers\_ensembl & 479\,bp
  & $\mathbf{0.905 \pm 0.012}$ & 0.706 & $+$0.199 \\
human\_enhancers\_cohn    & 500\,bp
  & $\mathbf{0.775 \pm 0.002}$ & 0.711 & $+$0.064 \\
human\_ocr\_ensembl       & 330\,bp
  & $\mathbf{0.811 \pm 0.009}$ & 0.806 & $+$0.005 \\
\bottomrule
\end{tabular}
\end{table}

\begin{figure}[h!]
\centering
\includegraphics[width=\linewidth]{fig4_genomic}
\caption{
  \textbf{Genomic sequence modeling results.}
  \textbf{(A)}~AUC on four genomic\_benchmarks (mean~$\pm$~std, three
  random seeds). FDM shows clear gains on promoter ($+$0.089) and
  enhancer tasks ($+$0.064 to $+$0.199); OCR performance is
  approximately competitive (near parity, $+$0.005).
  The larger variance on \texttt{human\_enhancers\_ensembl}
  (std~$= 0.012$) warrants caution.
  \textbf{(B)}~Downstream transfer.
  \emph{Top}: variant effect prediction on ClinVar (5-fold CV);
  fine-tuned FDM (AUC~$= 0.861 \pm 0.011$) approaches CADD (0.869),
  which uses manually curated features and cross-species conservation.
  \emph{Bottom}: cross-species transfer to \emph{Drosophila} enhancer
  recognition; human-pretrained FDM (0.792) outperforms a
  from-scratch CNN (0.736).
}
\label{fig:genomics}
\end{figure}

FDM shows clear gains on promoter and enhancer tasks (AUC $+$0.089
and $+$0.064 to $+$0.199 respectively), while remaining approximately
competitive on OCR ($+$0.005, near parity).
The larger variance on \texttt{human\_enhancers\_ensembl}
(std $= 0.012$) warrants caution.
Cross-chromosome generalisation (chr1--3 $\to$ chr19) yields
FDM AUC $= 0.903$ vs.\ CNN $= 0.895$, confirming that results do
not arise from chromosome-level memorisation.

\subsection{Variant Effect Prediction}

\begin{table}[h]
\centering
\caption{Variant effect prediction AUC (ClinVar chr17, 5-fold CV).
CADD\citenum{19} uses manually curated features and cross-species
conservation.}
\label{tab:vep}
\begin{tabular}{lc}
\toprule
Method & AUC \\
\midrule
FDM zero-shot            & 0.51 \\
CADD (reference)         & 0.869 \\
FDM fine-tuned (5-fold)  & $\mathbf{0.861 \pm 0.011}$ \\
\bottomrule
\end{tabular}
\end{table}

Fine-tuned FDM reaches AUC $= 0.861 \pm 0.011$, substantially above
zero-shot performance (0.51) and approaching CADD (0.869).
CADD incorporates manually curated genomic features and cross-species
conservation scores; direct comparison should be interpreted accordingly.

\subsection{Cross-Species Transfer}

\begin{table}[h]
\centering
\caption{Cross-species transfer AUC (\emph{Drosophila} enhancers\citenum{20},
$n = 6{,}914$).}
\label{tab:cross}
\begin{tabular}{lcc}
\toprule
Model & Species & AUC \\
\midrule
CNN (from scratch)           & \emph{Drosophila} & 0.736 \\
FDM-HG38 (human pretrained)  & \emph{Drosophila} & \textbf{0.792} \\
\bottomrule
\end{tabular}
\end{table}

Human-pretrained FDM achieves AUC $= 0.792$ vs.\ 0.736 for a
from-scratch CNN ($+$0.057), consistent with partial transfer of
regulatory sequence representations across species.
A preliminary mouse (mm10) experiment yielded FDM $= 0.741$
vs.\ CNN $= 0.718$ (577 samples; statistical power is limited).

%------------------------------------------------------------------
\section{Discussion}
\label{sec:discussion}

\paragraph{On the gradient sink and phase-aware training.}
The 7.5$\times$ improvement from Phase~0 alone ($487.5 \to 64.9$)
indicates the cache was severely undertrained, not that the architecture
was fundamentally limited.
Phase~1 provides further improvement ($64.9 \to 33.75$), confirming
complementary rather than redundant phases.
The gradient-ratio criterion provides an automatic trigger; its
robustness is empirical and would benefit from theoretical analysis.

\paragraph{On wave--cache functional divergence.}
MI divergence (wave: 0.474\,nats, cache: $\sim$0) and cache locality
bias are consistent with the intended decomposition but do not constitute
proof of strict functional separation.
Uniform $\theta$ across layers ($\lambda \approx 12.5$ tokens) indicates
operation at a single temporal scale; hierarchical behaviour would
require per-layer $\theta$ diversity, which we leave to future work.

\paragraph{On the OCR result.}
The near-parity on OCR (FDM: 0.811, HyenaDNA: 0.806) contrasts with
larger gains elsewhere.
We do not attribute this to architectural limitations with confidence;
establishing a causal link requires controlled ablations not performed here.

\paragraph{On the genomic--language gap.}
FDM achieves PPL $= 33.75$ on WikiText-103 and PPL $= 3.23$ on hg38,
suggesting the architecture transfers across sequence modalities.
Competitive results at scale would require substantially larger corpora.

\paragraph{Limitations.}
(1)~$\alpha$ adds a hyperparameter; behaviour beyond 137M parameters is
untested.
(2)~$K = 32$ global slots constrain long-range retrieval.
(3)~\texttt{enhancers\_ensembl} seed variance (std $= 0.012$) warrants
further investigation.
(4)~Cross-species results are preliminary.

\paragraph{Architectural directions.}
Uniform $W$ across layers and fixed $K$ are simplifying choices.
Layer-wise $W$ increasing with depth---shallow layers for local motifs,
deeper layers for longer context---is analogous to expanding receptive
fields\citenum{21}.
$K$ scaling with sequence length or task signal density could improve
long-range utility.
Both modifications require pretraining from scratch and are left to
future work.

%------------------------------------------------------------------
\section{Conclusion}

FDM decouples global state propagation from sparse local retrieval.
The structural tension between norm-preserving propagation and selective
retention motivates the architecture; the gradient sink problem motivates
phase-aware training.
On WikiText-103, FDM achieves perplexity 33.75---below our Transformer
reference---with measured runtime memory saturating at moderate context
lengths.
FDM shows clear gains on promoter and enhancer benchmarks and remains
competitive on OCR.
Fine-tuning on variant effect prediction approaches a feature-engineered
baseline.
These results suggest that explicit functional decomposition, combined
with training strategies that respect it, can yield competitive sequence
models with more functionally interpretable channel specialisation.

%------------------------------------------------------------------
\section*{References}

\begin{enumerate}\setlength{\itemsep}{1pt}\small

\item Merity, S.\ et al.\ Pointer sentinel mixture models.
      \textit{ICLR} (2017).

\item Nguyen, E.\ et al.\ HyenaDNA: Long-range genomic sequence modeling
      at single nucleotide resolution.
      \textit{NeurIPS} (2023).

\item Vaswani, A.\ et al.\ Attention is all you need.
      \textit{NeurIPS} (2017).

\item Gu, A., Goel, K.\ \& R\'{e}, C.\ Efficiently modeling long sequences
      with structured state spaces. \textit{ICLR} (2022).

\item Gu, A.\ \& Dao, T.\ Mamba: Linear-time sequence modeling with
      selective state spaces. \textit{arXiv}:2312.00752 (2023).

\item Peng, B.\ et al.\ RWKV: Reinventing RNNs for the Transformer era.
      \textit{EMNLP} (2023).

\item Dao, T.\ \& Gu, A.\ Transformers are SSMs: Generalized models and
      efficient algorithms. \textit{ICML} (2024).

\item Sun, Y.\ et al.\ Retentive network: A successor to Transformer.
      \textit{arXiv}:2307.08621 (2023).

\item Zhang, Z.\ et al.\ H2O: Heavy-hitter oracle for efficient
      generative inference. \textit{NeurIPS} (2023).

\item Xiao, G.\ et al.\ Efficient streaming language models with
      attention sinks. \textit{ICLR} (2024).

\item Beltagy, I., Peters, M.\ E.\ \& Cohan, A.\ Longformer: The
      long-document Transformer. \textit{arXiv}:2004.05150 (2020).

\item Zaheer, M.\ et al.\ Big Bird: Transformers for longer sequences.
      \textit{NeurIPS} (2020).

\item Dalla-Torre, H.\ et al.\ The Nucleotide Transformer.
      \textit{bioRxiv} (2023).

\item Zhou, Z.\ et al.\ DNABERT-2: Efficient foundation model for
      multi-species genomes. \textit{arXiv}:2306.15006 (2023).

\item Nguyen, E.\ et al.\ Sequence modeling and design from molecular
      to genome scale with Evo. \textit{Science} (2024).

\item Bengio, Y.\ et al.\ Curriculum learning. \textit{ICML} (2009).

\item Olsson, C.\ et al.\ In-context learning and induction heads.
      \textit{Transformer Circuits Thread} (2022).

\item Gre\v{s}ov\'{a}, K.\ et al.\ Genomic benchmarks.
      \textit{BMC Genomic Data} (2023).

\item Rentzsch, P.\ et al.\ CADD: Predicting deleteriousness of variants
      throughout the human genome.
      \textit{Nucleic Acids Research} (2019).

\item Stark, A.\ et al.\ Discovery of functional elements in 12
      \textit{Drosophila} genomes.
      \textit{Nature} \textbf{450}, 219--232 (2007).

\item He, K.\ et al.\ Deep residual learning for image recognition.
      \textit{CVPR} (2016).

\end{enumerate}

%------------------------------------------------------------------
\appendix
\section*{Supplementary Information}

\subsection*{S1.\ Query-Conditioned Wave State Modulation}

\emph{Auxiliary extension; not part of primary claims.}

\begin{equation}
  h_\mathrm{mod} = h_t \odot \bigl(1 + \tanh(W_\mathrm{ref} x_t)\bigr).
\end{equation}
Zero-initialisation of $W_\mathrm{ref}$ ensures $h_\mathrm{mod} = h_t$
at initialisation.
Multi-head orthogonal variant:
\begin{equation}
  h_\mathrm{mod} = h_t \odot \frac{1}{H}
    \sum_{i=1}^{H} \bigl(1 + \tanh(W_{\mathrm{ref},i} x_t)\bigr),\quad
  \mathcal{L}_\mathrm{orth} = \lambda \sum_{i \ne j}
    \|W_{\mathrm{ref},i} W_{\mathrm{ref},j}^\top\|_F^2.
\end{equation}
Modulation benefit concentrates in layer~0 ($\Delta\mathrm{PPL} = -1.36$),
diminishing rapidly in deeper layers.

\begin{table}[h]
\centering
\caption{Query-conditioned modulation ablation (starting PPL $= 64.9$).}
\label{tab:modulation}
\begin{tabular}{lc}
\toprule
Configuration & $\Delta$PPL \\
\midrule
Single-head, layer~0                & $-1.82$ \\
4-head orthogonal, layer~0          & $-2.13$ \\
8-head orthogonal, layer~0          & $\approx -2.1$ \\
Sequential layerwise (layers~0--11) & $-1.58$ \\
\bottomrule
\end{tabular}
\end{table}

\subsection*{S2.\ MQAR Associative Recall (Auxiliary)}

The following is a synthetic auxiliary experiment included for completeness;
it is not part of the primary claims.
Easy setting (sequence length~64, 8~KV pairs, vocabulary~512):
FDM $K=16$ accuracy $= 0.966$; Transformer $= 0.606$;
wave-only ($K=0$) $= 0.011$.
The near-zero wave-only score directly demonstrates that the cache
component is necessary for precise associative recall under this task,
complementing the mutual information analysis in the main text.



\begin{figure}[h]
\centering
\includegraphics[width=0.6\linewidth]{supp_memory}
\caption{
  \textbf{Runtime memory saturation of FDM under increasing context
  length.}
  FDM measured incremental memory remains approximately constant above
  $N = 2{,}048$ (396.9\,MB), whereas the Transformer KV cache grows
  linearly with sequence length (theoretical).
  At $N = 65{,}536$, FDM requires $4.6\times$ less incremental memory.
}
\label{fig:memory}
\end{figure}

\subsection*{S3.\ Training Details}

\textbf{Language model.}
$d = 576$, 12~layers, $n_G = 4$, $K = 16$, $W = 256$.
AdamW ($\beta_1=0.9$, $\beta_2=0.999$, wd$=0.01$).
Phase~0: lr$=10^{-4}$, steps~0--23{,}000.
Phase~1: lr$=3\times10^{-5}$, steps~23{,}000--40{,}000.
WikiText-103 (121M tokens). Hardware: single NVIDIA RTX~5090.

\textbf{Genomic model.}
$d = 512$, 8~layers, $n_G = 4$, $K = 32$, $W = 128$.
AdamW, lr$=10^{-4}$, cosine decay, 95{,}000~steps.
Fine-tuning: lr$=5\times10^{-5}$, patience~4, first~4 layers frozen.

\subsection*{S4.\ Genomic Benchmark Details}

Tasks from \texttt{genomic\_benchmarks}\citenum{18}, protocol follows
HyenaDNA\citenum{2}.
All results use checkpoint at step~54{,}000 (PPL~$=3.37$).
Seeds: 42, 123, 456.

\subsection*{S5.\ Cross-Species Experimental Details}

\textit{Drosophila}: positive from \texttt{drosophila\_enhancers\_stark}
\citenum{20} (chr2L/2R/3L/3R/chrX); balanced sampling
(3{,}457$+$3{,}457).
CNN baseline: 3-layer Conv1d (kernel~9) $+$ GAP $+$ 2-layer MLP.
Mouse (mm10): 577~balanced samples after N-masking.

\end{document}
